\chapter{Durchführung und Umsetzung}

\section{Projektstruktur}

Die Umsetzung ist als Python-Projekt mit PDM organisiert. Die wichtigsten Dateien und Ordner
sind:

\begin{itemize}
    \item \texttt{calib/mono\_then\_stereo\_calibration.ipynb}: Mono- und Stereo-Kalibrierung,
    \item \texttt{scene\_pointcloud\_from\_stereo.ipynb}: Aufnahme, Detektion, Rektifizierung,
          Disparität, Punktwolkenexport und Sanitation,
    \item \texttt{calib/mono\_camera\_params.pkl}: gespeicherte Mono-Intrinsiken,
    \item \texttt{calib/stereo\_params\_mono\_first.pkl}: gespeicherte Stereo-Parameter,
    \item \texttt{captures/}: gespeicherte linke und rechte Stereo-Rohbilder,
    \item \texttt{detections/}: YOLO-Ausgaben auf Rohbildern,
    \item \texttt{pointclouds/}: exportierte PLY-Punktwolken,
    \item \texttt{CustomModel/TrainingData/Stifterkennung.yolov11/}: YOLO-Datensatz,
    \item \texttt{best.pt}: Trainingsergebnisse.
\end{itemize}

Die wesentlichen Bibliotheken sind OpenCV, pypylon, Open3D, NumPy, Matplotlib und Ultralytics
YOLO. Der Quellcode ist notebookbasiert, damit Kalibrierung, Sichtprüfung und Export
schrittweise nachvollzogen werden können.


\section{Teil A: Bildverarbeitung}

\subsection{Kalibrierworkflow}

Die Kalibrierung folgt einem Mono-First-Ansatz. Zuerst wurden linke und rechte Kamera separat
kalibriert. Danach wurde die Stereo-Beziehung mit fixierten Mono-Intrinsiken bestimmt. Diese
Trennung erleichtert die Fehlersuche, weil schlechte Einzelkameraparameter nicht erst in der
Stereo-Kalibrierung sichtbar werden.

Die Konfiguration des Kalibriernotebooks umfasst ein ChArUco-Board mit $11 \times 8$
Quadraten, $10\,\mathrm{mm}$ Quadratgröße, $7\,\mathrm{mm}$ Markergröße,
\texttt{DICT\_4X4\_50} und aktivierter Legacy-ChArUco-Anordnung. Für gültige Bilder werden
mindestens zwölf interpolierte ChArUco-Ecken gefordert. Kameraordnung, Belichtung, Gain,
Fokus, Auflösung und Boardgeometrie müssen während einer Kalibriersitzung unverändert
bleiben.

\begin{table}[htbp]
    \centering
    \caption{Wesentliche Kalibrierparameter des aktuellen Stands.}
    \label{tab:calibration_values}
    \begin{tabular}{ll}
        \hline
        Parameter & Wert \\
        \hline
        Bildgröße & $1280 \times 1024\,\mathrm{px}$ \\
        Mono-RMS links/rechts & ca. $0{,}23\,\mathrm{px}$ / $0{,}24\,\mathrm{px}$ \\
        Stereo-RMS & ca. $1{,}34\,\mathrm{px}$ \\
        Baseline & ca. $123\,\mathrm{mm}$ \\
        Kalibriereinheit & Millimeter \\
        Ausgabedatei & \texttt{calib/stereo\_params\_mono\_first.pkl} \\
        \hline
    \end{tabular}
\end{table}

Die gespeicherten Stereo-Parameter enthalten die Kameramatrizen, Verzeichnungskoeffizienten,
Rotation, Translation, Essential- und Fundamental-Matrix, Rektifizierungsmatrizen,
Projektionsmatrizen und die Reprojektionsmatrix $Q$.

\begin{figure}[htbp]
    \centering
    \includegraphics[width=0.45\columnwidth]{Aufbau.jpg}
    \caption{Aufbau der Stereo-Kamera mit ChArUco-Board zur Kalibrierung.}
    \label{fig:stereo_calib_pair}
\end{figure}

\newpage


\subsection{Aufnahme und Rektifizierung}

Das Rekonstruktionsnotebook kann entweder neue Bilder mit den Basler-Kameras aufnehmen oder
das neueste gespeicherte Stereo-Paar aus \texttt{captures/} verwenden. Im Workspace liegen 60
Szenen- und Testbilder, also 30 linke/rechte Paare, vor. Vor dem Stereo-Matching kann CLAHE
aktiviert werden, um lokale Kontraste zu erhöhen.

\begin{figure}[htbp]
    \centering
    \includegraphics[width=0.47\columnwidth]{stereo_scene_20260615_170755_L.png}
    \includegraphics[width=0.47\columnwidth]{stereo_scene_20260615_170755_R.png}
    \caption{Beispiel eines gespeicherten Stereo-Rohbildpaars aus dem Ordner \texttt{captures/}.}
    \label{fig:stereo_raw_pair}
\end{figure}

Die Rektifizierung wird mit den gespeicherten OpenCV-Parametern durchgeführt. Dadurch liegen
korrespondierende Punkte idealerweise auf gleicher Bildhöhe. Dieser Schritt ist die
Voraussetzung dafür, dass StereoSGBM sinnvoll arbeiten kann.

\subsection{Disparität und Punktwolkenexport}

Für die Disparitätsberechnung wird \texttt{cv2.StereoSGBM\_create} eingesetzt. Die aktuelle
Konfiguration nutzt \texttt{minDisparity = 0}, initial \texttt{numDisparities = 384},
maximal \texttt{512}, eine Blockgröße von 7, \texttt{uniquenessRatio = 8},
\texttt{speckleWindowSize = 80}, \texttt{speckleRange = 2} und den Modus
\texttt{STEREO\_SGBM\_MODE\_SGBM\_3WAY}. Wenn mehr als fünf Prozent der gültigen Pixel am
oberen Disparitätslimit liegen, erweitert das Notebook den Disparitätsbereich automatisch.

Nach der Disparitätsschätzung werden nur Punkte im Tiefenbereich von
$400\,\mathrm{mm}$ bis $650\,\mathrm{mm}$ übernommen. Die Punktwolken werden als PLY-Dateien
gespeichert. Um sehr große Dateien handhabbar zu halten, wird beim Speichern auf maximal
400\,000 Punkte reduziert.

\begin{table}[htbp]
    \centering
    \caption{Erzeugte Bildverarbeitungsartefakte im Workspace.}
    \label{tab:vision_artifacts}
    \begin{tabular}{ll}
        \hline
        Artefakt & Anzahl / Zweck \\
        \hline
        \texttt{captures/} & 60 Rohbilder \\
        \texttt{detections/} & 26 annotierte YOLO-Bilder \\
        \texttt{pointclouds/} & 21 PLY-Punktwolken \\
        \texttt{calib/*.pkl} & gespeicherte Kalibrierparameter \\
        \hline
    \end{tabular}
\end{table}

\subsection{Datensanitation}

Zusätzlich zur rohen Punktwolke erzeugt das Notebook bereinigte Varianten. Der erste
Bereinigungsschritt schneidet die Disparitätskarte anhand definierter Disparitäts- und
Bildbereichsgrenzen. Der zweite Schritt arbeitet auf der gültigen Disparitätsmaske und
behält den größten zusammenhängenden Bereich. Kleine getrennte Komponenten werden ab einer
Mindestfläche von 2\,000 Pixeln gefiltert.

Beim letzten dokumentierten Lauf wurden nach der Disparitäts- und Tiefenfilterung 278\,810
Kandidatenpunkte betrachtet. Die Sanitation behielt 268\,574 Punkte, also 96,33 Prozent. Der
resultierende Bereich lag bei etwa $X=-77{,}1$ bis $197{,}3\,\mathrm{mm}$ und
$Z=523{,}9$ bis $580{,}7\,\mathrm{mm}$. Die bereinigte Datei
\texttt{scene\_pointcloud\_sanitized\_20260616\_091546.ply} hat eine Größe von ca.
$3{,}84\,\mathrm{MB}$.

\section{Teil B: Machine Learning}

\subsection{Datensatz}

Für die Stifterkennung wurde ein eigener YOLO-Datensatz erstellt. Die Ausgangsbilder liegen
zusätzlich nach Klassen sortiert in \texttt{CustomModel/TrainingData/}: 50 Bilder für
\textit{Marker}, 50 für \textit{Fineliner}, 50 für \textit{Pencil} und 60 gemischte Bilder.
Der daraus exportierte Roboflow/YOLO-Datensatz enthält drei Klassen:

\begin{itemize}
    \item Klasse 0: \textit{Fineliner},
    \item Klasse 1: \textit{Marker},
    \item Klasse 2: \textit{Pencil}.
\end{itemize}

Die Aufteilung im YOLO-Datensatz besteht aus 144 Trainingsbildern, 41 Validierungsbildern und
21 Testbildern. Die Labels liegen im YOLO-Format als normalisierte Bounding Boxes vor.

\begin{figure}[htbp]
    \centering
    \includegraphics[width=0.82\columnwidth]{labels.jpg}
    \caption{Label- und Bounding-Box-Verteilung des YOLO-Trainingsdatensatzes.}
    \label{fig:yolo_labels}
\end{figure}

\subsection{Training}

Trainiert wurde mit Ultralytics YOLO11m. Als Startpunkt wurde ein vortrainiertes Modell
verwendet. Die Trainingskonfiguration nutzte 640 Pixel Eingabegröße, Batchgröße 8,
Validierung während des Trainings, deterministische Ausführung mit Seed 0 und automatische
Optimizer-Auswahl. Die Trainingskonfiguration war auf 300 Epochen mit Patience 50 ausgelegt;
der gespeicherte Lauf umfasst 129 Epochen.

Zur Regularisierung und besseren Generalisierung wurden Standard-Augmentationen verwendet:
HSV-Anpassung, Translation, Skalierung, horizontales Spiegeln, Mosaic, RandAugment und
Random Erasing. Diese Augmentationen passen zur Empfehlung, kleine industrielle
Bilddatensätze künstlich zu variieren, damit das Modell nicht nur die konkreten
Trainingsaufnahmen auswendig lernt.

\begin{table}[htbp]
    \centering
    \caption{YOLO-Trainingsergebnisse des Laufs \texttt{stifterkennung\_yolo11m}.}
    \label{tab:yolo_metrics}
    \begin{tabular}{lll}
        \hline
        Metrik & Bester Wert & Epoche \\
        \hline
        Precision & 0,996 & 117 \\
        Recall & 0,967 & 129 \\
        mAP@50 & 0,985 & 128 \\
        mAP@50:95 & 0,845 & 79 \\
        \hline
    \end{tabular}
\end{table}

\begin{figure}[htbp]
    \centering
    \includegraphics[width=0.95\columnwidth]{results.png}
    \caption{Trainingsverlauf mit Verlusten und Validierungsmetriken des YOLO11m-Modells.}
    \label{fig:yolo_results}
\end{figure}

\begin{figure}[htbp]
    \centering
    \includegraphics[width=0.82\columnwidth]{confusion_matrix_normalized.png}
    \caption{Normalisierte Confusion Matrix des trainierten YOLO-Modells.}
    \label{fig:yolo_confusion_matrix}
\end{figure}

\subsection{Inferenz im Stereo-Workflow}

Die YOLO-Inferenz wurde in das Rekonstruktionsnotebook integriert, bleibt aber von der
Punktwolkenberechnung getrennt. Das Modell wird aus \texttt{best.pt} geladen, die
Konfidenzschwelle beträgt 0,4 und die IoU-Schwelle 0,45. Die Erkennung wird auf linkes und
rechtes Rohbild angewendet. Die Ausgaben werden mit Bounding Boxes gespeichert.

\begin{figure}[htbp]
    \centering
    \includegraphics[width=0.47\columnwidth]{yolo_scene_20260616_091521_L.png}
    \includegraphics[width=0.47\columnwidth]{yolo_scene_20260616_091521_R.png}
    \caption{YOLO-Inferenz auf linken und rechten Rohbildern vor der Rektifizierung.}
    \label{fig:yolo_raw_detection}
\end{figure}

\section{Zusammenwirken der Module}

Der aktuelle Workflow nutzt YOLO zur semantischen Analyse der Rohbilder und StereoSGBM zur
geometrischen Rekonstruktion. Diese Reihenfolge ist praktisch, weil die Objekterkennung nicht
durch Rektifizierungsartefakte beeinflusst wird. Für eine spätere 3D-Objektauswertung
könnten die 2D-Bounding-Boxes auf die rektifizierten Bilder übertragen und anschließend mit
den gültigen 3D-Punkten innerhalb der jeweiligen Objektregion verknüpft werden.
